% Source mirror: `paper/tex/intro.tex`.
\section{算法描述} \label{sec:intro}

基于哈希的签名是后量子数字签名中最有前景的候选方向之一，也可能是其中最保守的一条路线。它的优势在于，其经典与量子安全强度主要建立在密码哈希函数理想化困难性之上\footnote{对称原语（包括哈希函数）的设计哲学是：它们应当只允许通用攻击；否则该设计通常会被视为存在缺陷。}，因此相较于其他候选方案，其安全性更易理解，也更易评估。

\paragraph{有状态签名。} Ralph Merkle 提出了建立在 Lamport 一次性签名（OTS）之上的基于哈希的签名方案 \cite{C:Merkle89a,Lamport79}。随后，对有状态签名的改进工作发展出了扩展 Merkle 签名方案（XMSS）~\cite{hulsing2018xmss} 和 Leighton-Micali 签名（LMS）~\cite{mcgrew2019leighton}，并分别被 NIST~\cite{cooper2020recommendation} 与 IETF 标准化。

\paragraph{无状态签名。} 在典型的有状态签名中，例如 Merkle 签名方案（MSS），签名者必须记录哪些 OTS 私钥已经使用过，以避免重复使用带来的安全问题。但在许多实际场景中，这样的状态维护并不总是可行。Goldreich 提出了无状态的基于哈希的签名构造~\cite{C:Goldreich86a,Goldreich04}，消除了本地状态的需求，但代价是签名长度极大。近年来，这一研究方向重新受到关注。通过引入超树结构，SPHINCS~\cite{EC:BHHLNP15} 给出了 Goldreich 风格无状态哈希签名的一个实用化实例。SPHINCS 进一步成为 NIST PQC 标准化过程中多个后续方案的基础，包括 Gravity-SPHINCS~\cite{RSA:AumEnd18}、SPHINCS-Simpira~\cite{EPRINT:GueMou17} 以及 \SPHINCSP\ \cite{CCS:BHKNRS19}。更近期的优化与分析工作包括 SPHINCS+C~\cite{SP:HKRY23}、针对超树结构的更优尺寸-时间折中~\cite{C:KhoKudWag25}、可变层高/链长框架扩展~\cite{EPRINT:QinSun25}，以及重新审视常和 Winternitz 一次性签名并将其应用于 SPHINCS+ 与 XMSS 的 SPHINCS-alpha 工作~\cite{optimalwots}。NIST 还发布了无状态哈希签名的最终联邦标准 FIPS 205（SLH-DSA）~\cite{NIST:FIPS205}。

% Source mirror: `paper/figure/sha256-simple-comp.tex`.
\ifnum\draft=0

\begin{table*}[ht]
    \centering
    \caption{三组已打包 Kylin 家族的哈希调用与尺寸指标。密钥生成、签名和验证以哈希调用代理计数表示；公钥、私钥和签名大小以字节计。}
    \label{tab:comp tiny}
    \begin{tabular}{l r r r r r r}%
    \toprule
      实例 & 密钥生成 & 签名 & 验证 & 公钥 & 私钥 & 签名 \\ \midrule
      Kylin-128 & 141055 & 2038881 & 4482 & 32 & 64 & 7128 \\
      Kylin-256 & 198271 & 2691224 & 15507 & 64 & 128 & 26376 \\
      Kylin-512 & 519295 & 10723949 & 36480 & 128 & 256 & 83016 \\
     \vspace{-1em} \\ \bottomrule
    \end{tabular}
\end{table*}

\else
\begin{table}[th]
    \centering
    \caption{三组已打包 Kylin 家族的哈希调用与尺寸指标。密钥生成、签名和验证以哈希调用代理计数表示；公钥、私钥和签名大小以字节计。}
    \label{tab:comp tiny}
    \begin{tabular}{l r r r r r r}%
    \toprule
      实例 & 密钥生成 & 签名 & 验证 & 公钥 & 私钥 & 签名 \\ \midrule
      Kylin-128 & 141055 & 2038881 & 4482 & 32 & 64 & 7128 \\
      Kylin-256 & 198271 & 2691224 & 15507 & 64 & 128 & 26376 \\
      Kylin-512 & 519295 & 10723949 & 36480 & 128 & 256 & 83016 \\
     \vspace{-1em} \\ \bottomrule
    \end{tabular}
\end{table}

\fi


\paragraph{我们的工作概述。}
我们提出了 \SPHINCSA，一种无状态的基于哈希的签名方案。该方案在保留原始 \SPHINCSP\ 核心结构优势的同时，对这一最先进的无状态哈希签名方案进行了改进。

我们的优化主要来源于底层一次性签名方案的改进。我们证明，这一编码方案在所有树结构一次性签名中都具有尺寸最优性。进一步地，就我们目前所知，没有任何基于 DAG 的构造能够在尺寸效率上优于该方案，因此我们的方案在已知方案中也可视为尺寸最优~\cite{optimalwots}。

由于我们对 \SPHINCSP\ 的修改仅限于其底层 \WOTSP\ 一次性签名的编码方式，因此新方案的安全分析与原始方案基本一致。基于原始工作的安全分析，我们估算其安全级别，并仔细选择参数以获得更优的整体性能。

为了便于从提交包层面直接查看，\cref{tab:comp tiny} 给出了 \texttt{Kylin-128}、\texttt{Kylin-256} 与 \texttt{Kylin-512} 的哈希调用代理指标以及密钥/签名大小。具体而言，该表包含冻结参数集下的密钥生成、签名和验证哈希调用指标。

更完整的参数选择比较见 \cref{sec:performance}。如 \cref{tab:kylin-package-bench} 所示，已打包的 Kylin 实例给出了当前仓库本地环境下的密钥生成、签名与验证性能概况。对于更重视验证时间的特定场景，还可以进一步将参数重新调整到更有利于快速验证的方向。

\paragraph{致谢。}

我们感谢 \SPHINCSP\ 团队在原始 \SPHINCSP\ 项目中的出色工作。当前代码与文档中有相当一部分内容均继承自他们值得称道的成果。
